\documentclass{beamer}
\usepackage[utf8]{inputenc}

\usetheme{Madrid}
\usecolortheme{default}
\usepackage{amsmath,amssymb,amsfonts,amsthm}
\usepackage{txfonts}
\usepackage{tkz-euclide}
\usepackage{listings}
\usepackage{adjustbox}
\usepackage{array}
\usepackage{tabularx}
\usepackage{gvv}
\usepackage{lmodern}
\usepackage{circuitikz}
\usepackage{tikz}
\usepackage{graphicx}

\setbeamertemplate{page number in head/foot}[totalframumber]

\usepackage{tcolorbox}
\tcbuselibrary{minted,breakable,xparse,skins}

\definecolor{bg}{gray}{0.95}
\lstset{
    language=C,
    basicstyle=\ttfamily\small,
    keywordstyle=\color{blue},
    stringstyle=\color{orange},
    commentstyle=\color{green!60!black},
    numbers=left,
    numberstyle=\tiny\color{gray},
    breaklines=true,
    showstringspaces=false,
}

%------------------------------------------------------------
\title{4.2.14}
\date{September 14,2025}
\author{EE25BTECH11011 - Navya}

\begin{document}

\frame{\titlepage}

\begin{frame}{Question}
Find the direction and normal vector for the line 
\begin{align}
x + y = 0
\end{align}
\end{frame}

\begin{frame}{Solution}
The equation of line is :
\begin{align}
x + y = 0
\end{align}

This equation can be expressed in terms of matrices\\
Let
\begin{align}
\vec{x} = \myvec{x \\ y}
\end{align}
\begin{align}
\vec{n^T} = \myvec{1 & 1}
\end{align}
\begin{align}
c = 0
\end{align}

The line equation can be written as:
\begin{align}
\vec{n^T}  \vec{x} = c
\end{align}

\end{frame}

\begin{frame}{Solution}
Where $\vec{n}$ is the normal vector of the given line

The direction vector of the line can be found by observing the normal vector.
\begin{align}
\vec{m} = \myvec{-1 \\ 1}
\end{align}


This is true because if the director vector is represented as 
\begin{align}
\vec{m}  = \myvec{1 \\ m}    
\end{align}
then the normal vector can be represented as 
\begin{align}
\vec{n} = \myvec{-m \\ 1}
\end{align}
\end{frame}

\begin{frame}{Solution}
This can be verified by the following equation:
\begin{align}
\vec{n^T}\vec{m} = 0
\end{align}

\begin{align}
\myvec{1 & 1}\myvec{-1 \\ 1} = 0
\end{align}\\



The normal vector of the line is $\vec{n} = \myvec{1 \\ 1}$\\
The director vector of the line is $\vec{m} = \myvec{-1 \\ 1}$\\

From the figure, it is clearly verified that the theoretical solution matches with the computational solution.\\
\end{frame}


\begin{frame}[fragile]{ Code (C)}
\begin{block}{Part 1}
\begin{lstlisting}[language=C, basicstyle=\ttfamily\small, keywordstyle=\color{blue}, frame=single]
#include <stdio.h>

/* Dot product of 2D vectors */
int dot_product(int a[2], int b[2]) {
    return a[0]*b[0] + a[1]*b[1];
}

/* Check if two vectors are orthogonal */
int is_orthogonal(int a[2], int b[2]) {
    return dot_product(a, b) == 0;
}

/* Equation of the line: x + y = 0  -> y = -x */
double line_equation(double x) {
    return -x;
}

\end{lstlisting}
\end{block}
\end{frame}

\begin{frame}[fragile]{ Code (C)}
\begin{block}{Part 2}
\begin{lstlisting}[language=C, basicstyle=\ttfamily\small, keywordstyle=\color{blue}, frame=single]
/* Example test */
int main(void) {
    int n[2] = {1, 1};    // normal vector
    int m[2] = {-1, 1};   // direction vector

    printf("Dot product of n and m = %d\n", dot_product(n, m));
    printf("Are n and m orthogonal? %s\n", is_orthogonal(n, m) ? "Yes" : "No");

    for (double X = -5; X <= 5; X += 2.5) {
        printf("For x = %.2f, y = %.2f (on line x+y=0)\n", X, line_equation(X));
    }
    return 0;
}
\end{lstlisting}
\end{block}
\end{frame}

\begin{frame}[fragile]{Python Code}
\begin{block}{Part 1}
\begin{lstlisting}[language=Python, basicstyle=\ttfamily\small, keywordstyle=\color{blue}, frame=single]
import ctypes
import numpy as np
import matplotlib.pyplot as plt

# Load compiled shared library
lib = ctypes.CDLL("./code5.dll")

# Define signatures
lib.dot_product.argtypes = [ctypes.POINTER(ctypes.c_int), ctypes.POINTER(ctypes.c_int)]
lib.dot_product.restype = ctypes.c_int

lib.is_orthogonal.argtypes = [ctypes.POINTER(ctypes.c_int), ctypes.POINTER(ctypes.c_int)]
lib.is_orthogonal.restype = ctypes.c_int
\end{lstlisting}
\end{block}
\end{frame}


\begin{frame}[fragile]{Python Code}
\begin{block}{Part 2}
\begin{lstlisting}[language=Python, basicstyle=\ttfamily\small, keywordstyle=\color{blue}, 
frame=single]
lib.line_equation.argtypes = [ctypes.c_double]
lib.line_equation.restype = ctypes.c_double

# --- Line: x + y = 0 ---
normal_vector = (ctypes.c_int * 2)(1, 1)     # n = (1, 1)
direction_vector = (ctypes.c_int * 2)(-1, 1) # m = (-1, 1)
vector_origin = np.array([0, 0])

# Dot product + orthogonality check
dp = lib.dot_product(normal_vector, direction_vector)
print(f"Dot product of n and m: {dp}")

if lib.is_orthogonal(normal_vector, direction_vector):
    print("The vectors are orthogonal (as expected).")
else:
    print("The vectors are NOT orthogonal.")
\end{lstlisting}
\end{block}
\end{frame}

\begin{frame}[fragile]{Python Code}
\begin{block}{Part 3}
\begin{lstlisting}[language=Python, basicstyle=\ttfamily\small, keywordstyle=\color{blue}, frame=single]
# Compute points on line y = -x
x_vals = np.linspace(-5, 5, 200)
y_vals = [lib.line_equation(float(x)) for x in x_vals]

# --- Plotting ---
plt.style.use('seaborn-v0_8-whitegrid')
plt.figure(figsize=(8, 8))

# Line
plt.plot(x_vals, y_vals, label='Line: x + y = 0', color='blue', zorder=1)

# Direction vector
plt.quiver(vector_origin[0], vector_origin[1],
           direction_vector[0], direction_vector[1],
           angles='xy', scale_units='xy', scale=1,
           color='green', label='Direction Vector', zorder=2)
\end{lstlisting}
\end{block}
\end{frame}

\begin{frame}[fragile]{Python Code: Plotting}
\begin{block}{Part 4}
\begin{lstlisting}[language=Python, basicstyle=\ttfamily\small, keywordstyle=\color{blue}, frame=single]
# Normal vector
plt.quiver(vector_origin[0], vector_origin[1],
           normal_vector[0], normal_vector[1],
           angles='xy', scale_units='xy', scale=1,
           color='red', label='Normal Vector', zorder=2)

# Origin point
plt.plot(vector_origin[0], vector_origin[1], 'o', color='purple', markersize=8,
         label='Origin (0, 0)')

plt.title('Line x + y = 0 with Direction and Normal Vectors')
plt.xlabel('x-axis')
plt.ylabel('y-axis')

\end{lstlisting}
\end{block}
\end{frame}

\begin{frame}[fragile]{Python Code: Plotting}
\begin{block}{Part 5}
\begin{lstlisting}[language=Python, basicstyle=\ttfamily\small, keywordstyle=\color{blue}, frame=single]
plt.axis('equal')
plt.legend()
plt.grid(True)

plt.xlim(-5, 5)
plt.ylim(-5, 5)

# Save figure
plt.savefig("fig5.png")
plt.show()


\end{lstlisting}
\end{block}
\end{frame}

\begin{frame}{Graphical Representation}
    \begin{figure}
        \centering
        \includegraphics[width=1\columnwidth]{../beamer/figs/fig5.png}
        \caption{line x+y=0}
    \end{figure}
\end{frame}

\end{document}
